\documentclass{report}
\usepackage{graphicx}
\usepackage{fontspec}
\usepackage{fullpage}
\usepackage{polyglossia}
\usepackage{libertine}


\begin{document}
\pagestyle{empty}

\newcommand{\rowstyle}[1]{{\small \bfseries #1}}

\newcommand{\setin}[1]{%
 \clearpage
\begin{flushright}
  (#1) \hspace{\fill} \textbf{RNI No. MAHBIL /2009/40123}
\end{flushright}


\begin{center}
  \includegraphics[width=1.5cm]{indiaemblem}

  \vskip 40pt

\fontspec{#1}[Script=Devanagari, Mapping=devanagarinumerals]

{\Huge \bfseries महाराष्ट्र शासन राजपत्र}

  \vskip 20pt

{\Large \bfseries असाधारण भाग पाच-अ}





\rule[-\baselineskip]{\linewidth}{2pt}\\
\rule[4pt] {\linewidth}{0.5pt}\\ % Raising to 4pt centers the text vertically between the lines
  \rowstyle{वर्ष ३, अंक २०(१) ]} \hspace{\fill}
  \rowstyle{मंगळवार, ऑगस्ट ८, २०१७/ श्रावण १७, शके १९३९} 
  \hspace{\fill}
  \rowstyle{ [ पृष्ठे १८, किंमत: रुपये  ३६.00}  \\   % ₹ 
\rule{\linewidth}{0.5pt}\\
\rule[\baselineskip]{\linewidth}{2pt}



\textbf{असाधारण क्रमांक ४७३}

  \vskip 20pt

{\large\bfseries प्राधिकृत प्रकाशन }

  \vskip 20pt

{\bfseries महाराष्ट्र विधानसभेत  व महाराष्ट्र विधानपरिषदेत सादर केलेली विधेयके
  (इंग्रिजी अनुवाद) }


\end{center}

  \vskip 40pt

\begin{center}
{\large MAHARASHTRA LEGISLATURE SECRETARIAT}
\end{center}

The following Bill was introduced in the Maharashtra Legislative Assembly on the
8th August 2017 is published under Rule 117 of the Maharashtra Legislative Assembly Rules :
}

\setin{Shobhika}
\setin{CDAC-GISTSurekh}
\setin{Murty Sanskrit}


\end{document}

%%% Local Variables:
%%% mode: latex
%%% TeX-master: t
%%% End:
